% SGA 3, 第 VI_B 篇, 第一轮英文重构的完整中文翻译.
% 范围: 命题 3.9、定理 3.10 与引理 3.10.1 全文.
% 底本: Polo--Gille Exp6B-13oct24.pdf, 本组件第 20--22 页
% (合卷阅读本第 361--363 页; 原页边页码 345--347, 止于命题 3.11 前).
% OCR 与 Reinhold 英文 Markdown 仅用作比对见证.

\medskip
\phantomsection\label{prop:sga3-VIB-3.9}%
\noindent\textbf{命题 3.9.---}
\emph{设 \(G\) 是局部有限呈示的 \(S\)-群.
假设 \(G^0\) 可表. 则典范态射}
\[
  i:G^0\longrightarrow G
\]
\emph{是开浸入; 此外, \(G^0\) 在 \(S\) 上拟紧.}

由于 \(G^0\) 是函子 \(G\) 的子函子, 态射
\(i\) 是单态射, 从而是根态射. 根据函子 \(G^0\) 的定义,
由定义 (EGA~IV\(_4\), 17.1.1) 立刻可验证 \(i\) 是形式平展的
(注意 \(G^0\) 是 \(i\) 在 \(G\) 中的像).
最后, 根据利用所表示的函子刻画局部有限呈示 \(S\)-概形的准则
(EGA~IV\(_3\), 8.14.2(c)), 再由 \SGARefLink{cor:sga3-VIB-3.8}{3.8},
由于 \(G\) 局部有限呈示, \(G^0\) 也在 \(S\) 上局部有限呈示.
因此 \(i\) 局部有限呈示 (参见 EGA~IV\(_1\), 1.4.3).
所以它是 平展根态射, 从而是开浸入 (EGA~IV\(_4\), 17.9.1).

最后一个断言由 \hyperref[cor:sga3-VIB-3.6]{3.6} 得到.

\medskip
\phantomsection\label{thm:sga3-VIB-3.10}%
\noindent\textbf{定理 3.10.---}
\emph{设 \(G\) 是 \(S\)-群. 下列条件等价:}
\begin{enumerate}
\item[\textup{(i)}] \emph{\(G\) 在单位截面的各点处在 \(S\) 上平滑.}
\item[\textup{(ii)}] \emph{\(G\) 在单位截面的各点处在 \(S\) 上平坦且局部有限呈示,
  并且对每个 \(s\in S\), \(G_s\) 在 \(\kappa(s)\) 上平滑.}
\item[\textup{\SGAInlineRefTarget{thm:sga3-VIB-3.10:item:iii}(iii)}] \emph{\(G\)
  有一个在 \(S\) 上平滑的开子群概形 \(G'\).}
\item[\textup{\SGAInlineRefTarget{thm:sga3-VIB-3.10:item:iv}(iv)}] \emph{\(G^0\)
  可由 \(G\) 的一个在 \(S\) 上平滑的开子概形表示.}
\end{enumerate}

显然有
\[
  \textup{(iv)}\Longrightarrow\textup{(iii)}
  \Longrightarrow\textup{(i)},
\]
并且由 \hyperref[cor:sga3-VIB-1.3.1]{1.3.1} 和 \hyperref[cor:sga3-VIB-2.4]{2.4},
\textup{\SGARefLink{cor:sga3-VIB-1.3.1:item:i}{(i)}} 同时蕴含
\textup{\SGARefLink{cor:sga3-VIB-1.3.1:item:ii}{(ii)}} 与 \textup{(iii)}.
此外, 由 EGA~IV\(_4\), 17.5.1, \textup{(ii)} 蕴含 \textup{(i)}.

还须证明 \textup{\SGARefLink{thm:sga3-VIB-3.10:item:iii}{(iii)}} 蕴含
\textup{\SGARefLink{thm:sga3-VIB-3.10:item:iv}{(iv)}}.
下面的 \hyperref[lem:sga3-VIB-3.10.1]{引理~3.10.1} 表明 \(G'\) 包含 \(\underline{G^0}\), 且
\[
  G'^0=G^0.
\]
所以只须证明 \(\underline{G^0}\) 在 \(G\) 中开; 因为我们已在
\hyperref[case:sga3-VIB-3.4]{3.4} 中看到, 此时 \(G^0\) 将由
\(G'\) 在开子集
\[
  \underline{G'^0}=\underline{G^0}.
\]
上诱导的平滑子群概形表示. 因此可以假设 \(G'=G\).

为证明 \(\underline{G^0}\) 是开的, 只须证明每个 \(s\in S\)
在 \(S\) 中都有一个邻域 \(T\), 使
\[
  \underline{G^0}\cap\pi^{-1}(T)
\]
在 \(\pi^{-1}(T)\) 中开. 设 \(s\in S\). 由于 \(G=G'\),
态射 \(\pi\) 局部有限呈示. 因此, 与
\SGARefLink{prop:sga3-VIB-3.7}{3.7} 的证明一样, 可以构造 \(S\) 中包含 \(s\)
的开子集 \(T\), 以及 \(G\) 中包含 \(\epsilon(s)\) 的开子集 \(W\),
使得诱导态射
\[
  \pi'':W\longrightarrow T
\]
有限呈示, 并且以 \(\epsilon\) 诱导的态射
\[
  \epsilon'':T\longrightarrow W
\]
为截面. 对每个 \(t\in T\),
\[
  W\cap\underline{G_t^0}
\]
于是是 \((\pi'')^{-1}(t)\) 中包含 \(\epsilon''(t)\) 的连通分支.
由于 \(\pi\) 平滑, \(\pi''\) 也平滑, 因此它是有限呈示的平滑态射.
由 EGA~IV\(_3\), 15.6.5, 当 \(t\in T\) 时这些
\(W\cap\underline{G_t^0}\) 的并集 \(W^0\) 在 \(W\) 中开.

另一方面, VI\(_A\), \SGARefLink{sga3:vi:proposition:0-5}{0.5}, 给出
\[
  W^0\cdot W^0=\underline{G^0}\cap\pi^{-1}(T),
\]
须证明它在 \(\pi^{-1}(T)\) 中开. 现在可以假设 \(T=S\);
于是只须证明 \(W^0\cdot W^0\) 在 \(G\) 中开.
由于 \(\pi\) 普遍开, \(\mu\) 也普遍开
(VI\(_A\), \SGARefLink{sga3:VIA:0.1}{0.1}).
因此, 由于 \(W^0\) 在 \(G\) 中开,
\[
  W^0\cdot W^0=\mu(W^0\times_S W^0).
\]
也是开的.

这一结果将在 \SGARefLink{cor:sga3-VIB-4.4}{4.4} 中得到加强.

\medskip
\phantomsection\label{lem:sga3-VIB-3.10.1}%
\noindent\textbf{引理 3.10.1.---}
\emph{设 \(G\) 是纤维局部有限型的 \(S\)-群.
则 \(G\) 的每个开子群概形 \(H\) 都包含 \(\underline{G^0}\), 且满足}
\[
  G^0=H^0.
\]

设 \(s\in S\), 并令
\[
  G'_s=H_s\cap\underline{G_s^0}.
\]
则 \(G'_s\) 是 \(\underline{G_s^0}\) 的开子集,
并且因 \(\underline{G_s^0}\) 不可约而在 \(\underline{G_s^0}\) 中稠密
(VI\(_A\), \hyperref[cor:sga3-VIB-2.4]{2.4}).
因此, 由 VI\(_A\), \SGARefLink{sga3:vi:proposition:0-5}{0.5}, 有
\[
  G'_s\cdot G'_s=\underline{G_s^0}
\]
这表明
\[
  \underline{G_s^0}=G'_s\cdot G'_s\subset H_s\cdot H_s=H_s.
\]
所以对每个 \(s\) 都有 \(\underline{G_s^0}=\underline{H_s^0}\), 从而
\[
  \underline{G^0}\subset H,\qquad G^0=H^0.
\]
% Translation statistics: source_lines=140; translated_lines=136; inline_math=94; display_math=15; pending_TODO=0.
% Review: complete component translation; math fragments, reference keys, environments and note counts checked against source; formula prose translated; compilation pending parent integration.
